% Subseção da F-BALD-2 (menos baldeações contra mais rápido). Dono: Rafael.
% Fontes: docs/plano-de-desenvolvimento/features/feature-F-BALD-2.md (Como funciona e
% Resultados); números de benchmarks/medir_comparacao_baldeacoes.py em 17/11/2026.
% Revisado pelo Rafael em 19/11; exceção ao congelamento de 15/11 aprovada pelo grupo
% (docs/GUIA.md §8).
\subsection{Menos baldeações contra mais rápido}\label{sec:F-BALD-2}

Para quem prefere trocar menos de ônibus, a página Baldeações mostra quanto tempo isso
custa: o \emph{tempo a mais} do trajeto de menos baldeações em relação ao mais rápido e as
\emph{baldeações evitadas}. O mínimo $k$ vem de $G_L$ (Seção~\ref{sec:F-BALD}); o tempo, de
$G_E$, como no trajeto mais rápido, na mesma faixa e com as esperas $H/2$ da
Seção~\ref{sec:F-ESPERA}. Cronometrar o trajeto que a busca em $G_L$ devolve inflaria o
custo, porque o desempate dela não olha o tempo; por isso o tempo comparado é o do
\emph{mais rápido entre os trajetos com até $k$ baldeações}.

Esse trajeto sai de um Dijkstra em $G_E$ cujo estado é o par (vértice, embarques feitos):
a aresta de embarque soma um, as demais não, e não há estado com mais de $k + 1$ embarques.
Um estado é descartado quando o vértice já foi fechado com menos embarques, pois esse outro
chegou antes e pode fazer tudo o que ele faria. O custo é
$O((k + 2)(V + E)\log((k + 2)V))$ no pior caso, com $V$ e $E$ de $G_E$; na prática cada
vértice fecha poucas vezes. O resultado coincidiu com o do NetworkX num grafo em camadas
montado à parte em todas as buscas dos testes e do feed real.

No feed real, no pico da manhã ($G_E$ com $V = 46.100$ e $E = 208.956$), comparamos os
mesmos 200 pares sorteados da Seção~\ref{sec:F-BALD}. A comparação levou $212{,}9$~ms por
par em média e $1.015{,}5$~ms no pior caso (Intel Core i5-1235U). O pior caso passa um pouco
de 1~s, o limite do RNF-05; esse limite é da consulta de trajeto mais rápido, e a comparação
faz essa consulta, a de menos baldeações e, quando preciso, a busca com limite. Em 160 pares o mais rápido
já tinha o mínimo de baldeações. Nos outros 40, evitar baldeações custou $1{,}5$~min na
mediana (3\,\% do tempo do mais rápido) e $14{,}5$~min no máximo; 36 evitaram uma
baldeação e 4 evitaram duas. Por exemplo, entre a Rua Curitiba e a Rua Gentil Portugal do
Brasil, a linha 3053 direta leva $45{,}6$~min, contra $44{,}2$~min do mais rápido com duas
baldeações.

\emph{Limitação:} a comparação mostra um ponto da troca entre tempo e baldeações, não todos
os trajetos não dominados. Como $G_E$ encadeia caminhadas e $G_L$ admite uma de até 300~m
por vez, o mais rápido pode fazer menos baldeações que o mínimo de $G_L$; nesse caso nada é
evitado.
